\documentclass[10pt]{article}
\usepackage[margin=.68in]{geometry}
\usepackage{booktabs,graphicx,hyperref,xcolor,enumitem,tabularx,longtable,microtype,fontspec}
\setmainfont{TeX Gyre Heros}
\hypersetup{colorlinks=true,urlcolor=blue!55!black,linkcolor=blue!55!black}
\setlist{nosep,leftmargin=1.35em}
\setlength{\parskip}{.38em}
\setlength{\parindent}{0pt}
\newcommand{\code}[1]{\texttt{#1}}
\newcommand{\outdir}{artifacts/reports/qwen-code-native-full-study-20260915}
\newcommand{\figdir}{\outdir/figures}
\title{Native Qwen Code SFT Study\\\large All 36 held-out TRExFitter configuration tasks}
\author{ATLAS $H\rightarrow\gamma\gamma$ coding-agent project}
\date{15 September 2026}

\begin{document}
\maketitle

\begin{abstract}
We evaluate whether supervised fine-tuning (SFT) helps Qwen models carry out
small but consequential TRExFitter configuration changes through the real Qwen
Code 0.23.4 agent harness.  Every reported checkpoint is run on all 36 unique
validation tasks, replacing the earlier six-task pilot.  Under unpinned backend
defaults, Qwen3.5-4B base passes 21/36 tasks (58.3\%) while its one-epoch
rank-16 LoRA checkpoint passes 15/36 (41.7\%).  With decoding fixed at
temperature zero, top-p 0.95, and seed 1, the comparison instead becomes
24/36 versus 26/36.  Every evaluated
checkpoint below 4B scores 0/36 because it does not call native tools.  Loss alone is thus
not a reliable measure of end-to-end agent quality.  At 27B, base reaches
28/36 and LoRA reaches the study-best 29/36.  The report separates edit,
validation, protocol, scope, and harness failures; studies temperature and
model scale; and documents every training run without GPT-model comparisons.
\end{abstract}

\section{Problem}
A physicist expresses an analysis decision in natural language---for example,
``request MINOS errors for \code{mu\_H}'' or ``restore the missing \code{Sample:
WH} block.''  The agent must locate the relevant setting in
\code{analysis.config}, make the smallest semantically correct change, preserve
every unrelated decision, execute the requested Python verifier, and describe
only the evidence it actually observed.  A syntactically plausible answer is
not enough: a wrong luminosity, bin boundary, sample reference, or fit switch
can silently change the physics result.

The current SFT-stage objective is deterministic static correctness.  Running
TRExFitter and comparing output histograms would establish a stronger form of
semantic equivalence, but that expensive executable reward belongs to the
subsequent RL stage.

\section{Dataset and held-out benchmark}
Canonical examples are model-neutral structured \code{messages} plus optional
tool schemas.  They do not store rendered chat markup, token IDs, or hidden
reasoning.  The Qwen tokenizer/chat template creates model-specific native
function-call tokens at training time.  The exact training snapshot contains
129 replay-approved training trajectories and 38 validation conversations.
The public dataset release is
\href{https://huggingface.co/datasets/cxyang-ucb/hyy-sft}{cxyang-ucb/hyy-sft}.
The exact local Parquet SHA-256 values are retained with the run artifacts:
\begin{verbatim}
train       133b08ca562964d6e4158bd7ced3d257924dd2705e64097a4a80b8dc92695545
validation  2fcd9dc0cbd93ff5347652180c352d85b2f586e59c507ef0187f768a1cba7a7e
\end{verbatim}
Checkpoint snapshot revisions and rendered-template preflights are retained
there as well.

The execution benchmark is a separate set of 36 unique validation tasks: four
tasks in each of nine categories---fit controls, input repairs, job/execution,
normalization parameters, references, region binning, region presentation,
sample settings, and whole-block reconstruction.  Splits are made by logical
task identity so a paraphrase of the same operation cannot leak across splits.
Each task starts from its own mutated configuration and has a machine-readable
change contract and evidence requirement.

\subsection{What one trajectory looks like}
A deterministic training demonstration asks for \code{FitBlind: TRUE}.  Its
assistant messages call \code{edit} with a minimal replacement, observe the
native tool result, call \code{run\_shell\_command} with the exact verifier,
observe exit code zero, and only then report success.  The canonical row stores
those structured messages and calls, not the rendered Qwen syntax.

One held-out task asks the agent to restore \code{UseMinos: mu\_H}, which tells
the statistical fit to compute asymmetric MINOS uncertainty for the Higgs
signal-strength parameter.  Both the 4B base and default LoRA checkpoint solve
this example: they read the file, replace the commented line, and obtain a
successful verifier result.  The scorer then independently confirms the
requested value and preservation of every unrelated setting.  This is one
illustrative success, not evidence for the aggregate comparison by itself.

\begin{verbatim}
%  UseMinos: mu_H        ->        UseMinos: mu_H
python -m trex_fitter.config_verify analysis.config --actions n \
  --evidence-level S
\end{verbatim}

\section{Native execution and independent scoring}
Qwen Code owns the agent loop and exposes its three bare-mode built-ins:
\code{read\_file}, \code{edit}, and \code{run\_shell\_command}.  The model can
use ordinary search commands and runs
\code{python -m trex\_fitter.config\_verify}; no repository-defined config-edit
tool is exposed.  Each task runs in an initialized, disposable Git workspace.
The repository scorer runs only after Qwen Code exits and cannot be called by
the model as a shortcut.
For the 27B runs, four disjoint workers share one four-GPU inference server to
fit the allocation.  This changes only scheduling: each of the 36 tasks still
has a separate Qwen Code process, Git workspace, event stream, and score, and
the combined task IDs are checked for uniqueness before reporting.

\textbf{Strict success} means that all of the following independently checked
conditions hold:
\begin{enumerate}
  \item \textbf{Requested edit}: the named setting has the requested semantic value;
  \item \textbf{Preservation}: every unrelated semantic setting is unchanged;
  \item \textbf{Static evidence}: parsing and the task's S/C/I evidence level pass;
  \item \textbf{Verifier protocol}: a successful exact verifier call is observed
        after the final edit;
  \item \textbf{Workspace scope}: successful native calls stay inside the
        disposable task and use only approved native tools;
  \item \textbf{Harness completion}: the Qwen Code process completes normally.
\end{enumerate}
The first three diagnose file correctness.  The last three diagnose whether an
agent completed the requested tool protocol safely.  Failures can overlap.
The verifier's one-letter evidence levels are: \textbf{S}, static syntax,
types, and references only; \textbf{I}, static checks plus inspection of input
ROOT-file metadata; and \textbf{C}, static checks plus compatibility with the
project's Coffea execution path.  Coffea is the Python event-processing system
used by this analysis.  The letter is a required verification depth, not a
model score.

\section{Training setup and reasoning behavior}
VERL performs assistant-only SFT with BF16, one epoch, seed 1, and no
truncation.  Runs through 9B use global batch 8 and 16 optimizer steps.  The
27B capacity run uses global batch 16 and eight steps across 16 A100s.  The
maximum sequence length is 16,384 through 9B.  The 27B tokenizer renders its
longest training row as 13,629
tokens, so that run uses 14,336 without truncation.  A 16,384-token attempt ran
out of memory even on eight A100s; a 12,288-token attempt correctly stopped on
the first row longer than its declared limit.  The final 27B configuration
uses four-way Ulysses sequence parallelism, which distributes token positions
within each group, plus removal of padding tokens before model computation.
Preflight
checks confirm correct assistant-span masks, identical turn-wise/full
rendering, the pinned tokenizer/template revision, and
\code{enable\_thinking=False}.  LoRA (low-rank adaptation) freezes the base
weights and learns small rank-16 updates; full SFT updates every model weight.
An assistant-span mask marks exactly which assistant-written tokens contribute
to the training loss; user prompts and tool outputs are context, not targets.
BF16 is a memory-efficient 16-bit number format.  An optimizer step is one
parameter update; the learning rate (LR) controls its size.  An A100 is the GPU
accelerator used for both training and serving.

\input{\outdir/training-results.tex}
\begin{table}[ht]
\centering\small
\begin{tabular}{@{}lcccccc@{}}
\toprule
Checkpoint & Update & LR & A100s & Steps & Held-out loss & Peak GB/GPU \\
\midrule
\trainingresultrows
\bottomrule
\end{tabular}
\caption{Fresh training on the Qwen Code-native snapshot.  Held-out loss is
assistant-token imitation loss, not execution accuracy.}
\end{table}

The Qwen3.5 checkpoints are reasoning-capable; the 0.5B comparison instead uses
Qwen2.5-Coder.  Hidden reasoning is not an SFT target.  For Qwen3.5, the serving
template's default enables thinking and the native event streams explicitly
contain \code{thinking} blocks.  Thus those models do use inference-time
reasoning, while training supervises only visible assistant/tool behavior.  A
controlled no-thinking evaluation tests this train/inference mismatch separately.

\begin{figure}[ht]
\centering\includegraphics[width=.72\textwidth]{\figdir/thinking-ablation.pdf}
\caption{The same 36 tasks with the serving template's thinking behavior on or
explicitly disabled.  This isolates inference-time reasoning from SFT.}
\end{figure}

Disabling thinking changes base performance modestly, from 21/36 to 19/36,
but collapses the default LoRA checkpoint from 15/36 to 1/36.  The disabled
runs contain zero thinking blocks, confirming that the control took effect.
This checkpoint therefore depends strongly on inference-time reasoning even
though hidden reasoning was excluded from its SFT loss.  The result exposes a
train/inference mismatch; it does not justify adding private chain-of-thought
to the dataset, because visible planning and recovery behavior can instead be
supervised and scored.

\begin{figure}[ht]
\centering\includegraphics[width=.9\textwidth]{\figdir/training-curves.pdf}
\caption{Training loss by optimizer step; diamonds are final held-out loss.
Different model sizes are not expected to share an absolute loss scale.}
\end{figure}

\subsection{LoRA hyperparameter ablation}
\begin{figure}[ht]
\centering\includegraphics[width=.9\textwidth]{\figdir/hyperparameter-ablation.pdf}
\caption{Three one-epoch 4B LoRA runs compare rank and learning rate.  The
left panel is native execution; the right panel is trajectory-imitation loss.
Both are needed because a lower loss need not produce a better agent.}
\end{figure}

The default run uses rank 16 and learning rate $10^{-4}$.  One ablation lowers
the learning rate tenfold, and another halves the adapter rank while preserving
the default learning rate.  Global batch, data, seed, context length, and epoch
count are fixed; execution uses temperature zero, top-p 0.95, and seed 1 for
all three checkpoints.  The default used four GPUs whereas the new ablations used one
each; the global optimizer batch is unchanged, but small distributed-numerical
differences mean this is not a mathematically exact one-variable experiment.
Rank 16 at LR $10^{-4}$ reaches held-out loss 0.079 and 26/36 strict successes.
Lowering LR to $10^{-5}$ underfits the demonstrations (loss 0.535) but still
reaches 24/36.  Halving the rank gives loss 0.089 and 27/36, one task above the
default.  On this sample, rank 8 is the more economical choice, but the one-task
accuracy difference is too small to claim a stable quality improvement.

\section{Checkpoint results}
\input{\outdir/checkpoint-results.tex}
\begin{table}[ht]
\centering\small
\begin{tabular}{@{}lccccc@{}}
\toprule
Checkpoint & Tasks & Strict success & CLI failures & Native calls & Thinking blocks \\
\midrule
\checkpointresultrows
\bottomrule
\end{tabular}
\caption{All entries use the same 36 tasks.  ``CLI failures'' means Qwen Code
returned nonzero or timed out; it is not automatically an editing failure.}
\end{table}

\begin{figure}[ht]
\centering\includegraphics[width=.94\textwidth]{\figdir/strict-success-by-checkpoint.pdf}
\caption{Strict end-to-end success.  Gray bars are unmodified base models;
blue bars are SFT checkpoints.  Error bars are 95\% Wilson binomial intervals;
that is a small-sample uncertainty range for the observed success fraction.
They are wide because the benchmark contains only 36 tasks.}
\end{figure}

\begin{figure}[ht]
\centering\includegraphics[width=.62\textwidth]{\figdir/controlled-4b-pair.pdf}
\caption{The causal 4B comparison with identical temperature zero, top-p 0.95,
seed 1, harness, prompt, and 36 tasks.  This removes the backend-default
sampling ambiguity in the original pair.}
\end{figure}

The controlled comparison gives base 24/36 and LoRA 26/36, a two-task LoRA
advantage.  This is directionally different from the backend-default pair and
shows that decoding is a material experimental variable.  With only 36 tasks,
neither small difference should be presented as a stable model ranking without
additional seeds.

Under backend-default sampling, the 4B result reverses what training loss alone
might suggest.  LoRA lowers
held-out imitation loss to 0.079, yet strict execution falls by six tasks from
base.  The SFT checkpoint makes more tool calls and emits more thinking blocks,
but activity is not correctness.  It also has more CLI failures.  At 0.8B,
base, LoRA, and full SFT all produce one thinking response per task and no tool
calls, so neither optimization method crosses the basic agent-capability
threshold.  The 0.8B full-versus-LoRA comparison therefore shows no execution
advantage for updating all weights despite the lower full-SFT held-out loss.
The 0.5B and 2B base/full pairs likewise make no native calls and score zero.
At 9B, base reaches 20/36 and LoRA reaches 19/36: the gap is smaller than at
4B, but SFT still does not improve strict success and increases CLI failures
from four to nine.  These are descriptive results on 36 tasks, not sufficiently
large samples for fine-grained statistical claims.  Scaling to 27B produces the
best absolute results: base reaches 28/36 with two CLI failures and LoRA reaches
29/36 with one.  Their category profiles differ---LoRA improves input repair,
job/execution, and region presentation, but loses one fit-control and one sample
setting task---so the one-task aggregate gain is not broad dominance.

\section{Where checkpoints fail}
\begin{figure}[ht]
\centering\includegraphics[width=.98\textwidth]{\figdir/diagnostic-gates.pdf}
\caption{A bar counts a gate independently.  ``Static evidence'' means the
final file passes the required deterministic checker; ``verifier observed''
means the agent itself successfully ran the exact command after its last edit.}
\end{figure}

\begin{figure}[ht]
\centering\includegraphics[width=.98\textwidth]{\figdir/category-results.pdf}
\caption{Category-level strict results.  Each cell has exactly four tasks, so
labels show counts rather than potentially misleading percentages.}
\end{figure}

The gate plot distinguishes three common stories.  First, a model can leave a
valid configuration but fail to make the requested change (edit-contract
failure).  Second, it can make the correct change but omit or fail the required
verifier (protocol failure).  Third, the CLI or endpoint can fail even when a
partial edit is reasonable (harness failure).  Treating all three as simply
``wrong answer'' would hide whether better demonstrations, decoding, or systems
engineering is needed.

\section{Temperature study}
\begin{figure}[ht]
\centering\includegraphics[width=.72\textwidth]{\figdir/temperature-study.pdf}
\caption{Each temperature is evaluated on all 36 tasks with seed 1 and
top-p 0.95.  Temperature controls randomness: zero selects the most likely
continuation, while larger values permit more varied choices.  Top-p 0.95
restricts sampling to tokens covering the most likely 95\% of probability;
the seed fixes the pseudo-random starting point.}
\end{figure}

The original matched runs used Qwen Code/SGLang model defaults.  In the new
sweep, Qwen Code remains in bare mode, so its native tool definitions,
function-call parser, execution loop, and messages are unchanged.  A transparent
localhost proxy modifies only the OpenAI-compatible request sent to the local
server, injecting temperature, top-p, and seed.  The run manifest records these
overrides, making the sampling condition auditable without replacing the agent
harness.

At seed 1, temperature zero passes 26/36 tasks; temperatures 0.2 and 0.7 each
pass 23/36.  This favors deterministic decoding for the present checkpoint,
but a single seed cannot characterize sampling variance at nonzero temperature.
The earlier backend-default result (15/36) is not placed on this controlled
curve because its effective sampling parameters and random seed were not pinned.

\section{Training throughput versus A100 count}
\begin{figure}[ht]
\centering\includegraphics[width=.72\textwidth]{\figdir/a100-throughput.pdf}
\caption{Fresh fixed-step VERL measurements on the same native SFT snapshot.
The metric is total non-padding training-sequence tokens divided by wall time;
the first warm-up step is excluded and the median resists checkpoint-I/O outliers.}
\end{figure}

Throughput is an efficiency measurement, not model quality.  Scaling below the
dashed line means communication, padding imbalance, or fixed overhead prevents
perfect linear speedup.  It helps choose the smallest useful GPU count: adding
GPUs is justified only when saved wall time outweighs queue and communication
costs.  The measured medians are 7,330, 12,768, and 22,508 sequence tokens/s on
one, two, and four A100s.  These correspond to 1.74$\times$ and 3.07$\times$
speedup, or 87\% and 77\% parallel efficiency relative to one GPU.  Four GPUs
therefore shorten a fixed run substantially, but do not deliver a fourfold
speedup; using four nodes for this small checkpoint would be inefficient.

\clearpage
\section{Conclusions}
The 36-task benchmark supports a more careful conclusion: the present SFT
recipe does not yield a decoding-robust improvement in native Qwen Code
performance.  It loses six tasks under backend defaults but gains two under the
controlled temperature-zero setting, despite similarly low validation loss.
At 27B, LoRA improves by only one task (28 to 29), again smaller than the
benchmark's sampling uncertainty.
The most likely limiting factors are the small number of demonstrations,
insufficient native Qwen recovery trajectories, distribution mismatch between
deterministic gold traces and real agent behavior, and optimization sensitivity.
The next SFT round should increase independently replayed native successes and
recoveries, retain the 36 tasks as untouched evaluation, and select checkpoints
by strict execution accuracy rather than loss alone.  VERL remains appropriate
because the same model/data stack can proceed to executable-reward RL after the
SFT baseline is sound.

\section*{Reproducibility}
Every run stores a manifest, the Qwen Code version and sampling controls, initial
and final configs, native JSON event streams, stderr, per-task scores, and a
36-row summary.  A rescore-only mode reapplies current independent gates without
rerunning a model.  An absolute path is accepted only when it resolves inside
the task workspace; other machine paths remain a scope failure.  The runner's
unit suite covers native call parsing, verifier ordering, semantic equivalence,
workspace scope, and transparent request-level sampling controls.

\appendix
\section{One failure example for every checkpoint and gate}
``No failure observed'' means that model passed that particular diagnostic gate
on all 36 tasks; it does not imply strict success on every task.
\scriptsize
\input{\outdir/failure-examples.tex}
\begin{longtable}{@{}p{.12\textwidth}p{.13\textwidth}p{.27\textwidth}p{.41\textwidth}@{}}
\toprule
Checkpoint & Failure type & Task & Observed example \\
\midrule
\endfirsthead
\toprule
Checkpoint & Failure type & Task & Observed example \\
\midrule
\endhead
\failureexamplerows
\bottomrule
\end{longtable}

\end{document}
